\chapter{Concluding remarks} \label{concluding}

The goal of this book was to present an analysis of the Japanese DP and its internal modifiers and to transfer the findings into cross-linguistic perspective. In this final chapter, I present the main findings, discuss some additional points and raise questions for future research.

\section{Main findings}

\subsection{Summary}

The following main findings emerged over the course of the discussion.

\ex. \a. The Japanese DP is made up of a DP-external and a DP-internal domain. The DP-internal domain is split up into a domain for direct and a domain for indirect modification. \is{left periphery}
\b. The DP-external domain, the left periphery, hosts functional projections for quantificational modifiers, non-restrictive relative clauses, potentially focused modifiers, as well as unspecific functional projections for (restrictive) indirect modifiers as well as direct modifiers. These modifiers move to this domain independently.
\c. Modifiers within the DP are licensed by an additional functional head, called \textsc{mod}, which licenses them as attributive modifiers. It licenses all modifiers, irrespective of word class and modification type. I argued that this head is universal. It is realized as /i/ and /na/ following the two adjective groups and as /no/ elsewhere. It is not realized when the tense head is overt.
\d. On the other hand, the complex elements following Japanese adjectives and \textit{no}-modifiers in indirect structures are combinations of the predicative copula, the dummy copula and a tense marker.
\e. Direct modifiers exist in Japanese. However, Japanese suggests that the direct domain consists of a possible infinite number of functional projections which are not necessarily dedicated to fine-grained semantic categories. This explains the lack of ordering principles observable in this language. The direct domain in languages such as English, on the other hand, does contain fine-grained functional projections. The reason for this can be sought in parametric differences. In other words, Japanese lacks the language-specific inventory (Units of Languages) to realize specific functional projections. There are also in-between languages, e.g. \is{Persian} Persian and \is{Modern Greek} Modern Greek, in which a subpart of the functional projections are realized, and it seems that if functional projections are realized, it is always those reserved for direct-only modifiers.
\f. Indirect modifiers only come in one size, they are CPs. The indirect domain is higher than the direct domain and encompasses several functional projections above and below a projection hosting numerals. Indirect modifiers are base-generated below this projection but can independently move to any projection above FP\textsubscript{Num}. More complex indirect modifiers are more likely to be subject to such movement operations.

I argued that the structure of the DP is uniform across languages, i.e. all languages share the same universal spine. Importantly, I left room for cross-linguistic variation as to how this spine is filled. Japanese presents (up to now) the most underspecified case in which functional projections are unspecific and not dedicated to modifiers of different types, word classes, semantic categories or sizes. Other languages are more precise in their structure of the DP, and for each relevant domain. In \figref{fig:Structure of the Japanese DP3}, I repeat the DP-structure, based on Japanese, proposed in this book.


\begin{figure}
\small
 \begin{subfigure}{0.5\textwidth}

 \begin{tikzpicture}[baseline=0pt, remember picture]
\tikzset{level distance=30pt, sibling distance=-5pt}
 \Tree [.DP\textsubscript{external} {} [.D′\textsubscript{external} [.FP\textsubscript{\textsc{quantity}} XP [.F\textsubscript{\textsc{quantity}} [.FocP XP [.Foc′ [.FP\textsubscript{n} XP [.F′\textsubscript{n} [.FP\textsubscript{n+1} XP [.F′\textsubscript{n+1} [.\textbf{DP\textsubscript{internal}} {} ... ] F ] ] F ] ] Foc ] ] F ] ] D ]]
\end{tikzpicture}
 \end{subfigure}%
 \begin{subfigure}{0.5\textwidth}
 \begin{tikzpicture}[baseline=0pt, remember picture]
\tikzset{level distance=30pt, sibling distance=-5pt}
 \Tree [.\textbf{DP\textsubscript{internal}} {} [.D′\textsubscript{internal} [.FP\textsubscript{indirect-n} XP [.F′\textsubscript{indirect-n} [.FP\textsubscript{Num} NumP [.F′\textsubscript{Num} [.FP\textsubscript{indirect-n+1} XP [.F′\textsubscript{indirect-n+1} [.\textbf{dP} [.FP\textsubscript{direct-n} XP [.F′\textsubscript{direct-n} [.FP\textsubscript{direct-n+1} XP [.F′\textsubscript{direct-n} [.$\sqrt{}$P XP [.$\sqrt{}$′ \edge[roof]; {N} ] ] F ] ] F ] ] d ] F ] ] F ] ] F ] ] D ] ]
\end{tikzpicture}
 \end{subfigure}
 \caption{Structure of the Japanese DP}
 \label{fig:Structure of the Japanese DP3}
\end{figure}

\subsection{What Cartography can offer for Japanese}

I hope to have shown in the preceding discussion that the cartographic research field provides the necessary tools and diagnostics for an analysis of the Japanese DP, and allows us to integrate Japanese into the by now large group of languages which have been subject to an in-depth analysis in this research program. Despite the fact that most existing analyses are dedicated to well-known language groups such as Germanic and Romance, it is astonishing how well these tools lend themselves for an application to other language groups. Most of the findings reported here are in line with the assumptions made in Cartography, which means Japanese provides further evidence for its validity. This concerns especially the inventory of modifiers, the dichotomy between direct and indirect modification and their connection to the different DP-domains.

\subsection{What Japanese can offer for Cartography} \label{discussion u20}

However, it is important to acknowledge that the application of Cartography to the Japanese DP is only possible if one makes the correct amendments and is open to draw critical conclusions, which brings me to the second important insight of this study: the contribution Japanese can offer to Cartography.

Perhaps the most important contribution is that the notorious lack of ordering restrictions, which has for decades been the main argument to deny direct (attributive) modification in Japanese, is not problematic at all; rather the presence or absence of ordering restrictions and type of modification are two orthogonal issues. Japanese suggests that all languages start out with unspecified functional projections and that only some languages can realize specific functional projections. The fine-grained hierarchy of functional projections is, on the one hand, one of the ground pillars of Cartography; on the other hand, it has constantly been the fuel for criticism put forward by various scholars (cf. \citealt{larson:2021:rethinking}).

By integrating the ideas that some languages exhibit more specific functional projections than others together with the concept of paradigmatic variability into the research field, I proposed that we can expand our knowledge of the DP and take this as a starting point for further investigation.

At this point, I would like to emphasize that the account proposed in this book, and built mostly on Japanese data, in no way undermines the core predictions and achievements of Cartography. This concerns, most of all, perhaps its biggest achievements: reducing left-right asymmetries and different orders of modifiers in relation to the lexical head, while at the same time explaining unattestable \citep{cinque:2025:superiority} orders through different movement operations the lexical head can undergo (see Section \ref{left-right}).\footnote{I am grateful to an anonymous reviewer for also raising this point.} \is{left-right asymmetries} \is{U20}

The system proposed here shares with standard Cartography the assumption of distinct merging domains for demonstratives, numerals and adjectives. In fact, Japanese presents a language where this order is reflected in linear order to a maximal extent. Assuming demonstratives are merged in DP (here: in Spec, DP), numerals within the indirect domain, and adjectives within the direct domain, if direct, or the indirect domain, if indirect, but then in a lower position than NumP, the relative order is unaffected by the broad direct and indirect domains. Applying this, for instance, to a language such as \is{Persian} Persian, which reflects the basis order Dem $\succ$ Num $\succ$ N $\succ$ A, the noun moves across each modifying adjective. This means that these general assumptions are unaffected by the idea of unspecific functional heads.

The only caveat exists with regard to (direct) adjectives, or more broadly \textit{modifiers}, where, as I argued, language variation comes into play. In Japanese, (direct) modifiers do not merge in a specific order, hence there is no fixed underlying order. However, since the noun does not move in this language anyway, this observation, only with regard to Japanese, does not interfere with U20-effects. Should a counterpart language to Japanese be found, and \is{Persian} Persian at least shows certain indications, then free order of adjectives (or: modifiers) in postnominal position will hold.

On the other hand, in languages that do have the language-specific entities to form specific functional projections in the direct domain, the relevant ordering principles outlined in Section \ref{left-right} will hold regardless. Specifically, only those orders are possible that are the result of the various movement operations accessible to the noun, and combinations thereof. Unattestable orders are ruled out by virtue of not being derivable by these operations. It is a fact that when we look at the variety of languages investigated through a cartographic lens, astounding parallels in terms of ordering emerge \citep[§ 2]{cinque:2025:superiority}, giving evidence that the principles of ordering and of realizing functional projections are uniform. \is{left-right asymmetries} \is{U20}

In other words, the core assumptions, and achievements, of Cartography are unaffected by the theory and the empirical data points presented throughout this book, however, certain assumptions, specifically the nature of the functional projections within the direct and the indirect domain and their strict mapping to order, are challenged by these aspects, which eventually has led to an evolution of cartographic theory.

Another crucial amendment, that enriches the DP structure, is the integration of the syntactic head \is{MOD@\textsc{mod}}\textsc{mod} as a universal mediator of modification as it offers a reference to the internal structure of modifiers. Tightening the connection between the structure of the DP and the internal structure of its modifiers can only be beneficial for future research, as it will allow us to draw conclusions from the macro- to the micro-perspective and vice versa.

\section{Open questions and directions for future research}

This study has probably opened as many new questions as it has answered, which I consider a good thing. Below, I will outline some immediate prospects, topics and concerns.

\subsection{Cartography: Where do we go from here?}

Assuming the DP structure I proposed here is universal and Japanese comes close to the baseline language in which almost no specific, or dedicated, functional heads exist within the direct, and indirect, domain begs the question whether similar languages exist. If they do, what could such languages have in common? Can we derive arial, genealogical or syntactic similarities and draw typological generalizations? In the general spirit of Cartography, precise typological considerations ensure the system is as theoretically and empirically sound as possible. A precise typology of languages according to the degree to which they adhere to ordering restrictions or, more importantly, allow ordering violations is therefore imperative for further progress in cartographic theory.

In addition to the rather lenient ordering observable in Japanese and the rather strict ordering observable in many other languages, what about further in-between languages besides the ones investigated here (Section \ref{language variation})?

The same applies for the indirect domain. Contrary to the situation in the direct domain, I expect that we find a large variety of languages of the Japanese type, given the fact that many (not only prenominal relative clause) languages seem to abide by the same principles, namely to only show one type of relative clause overtly. Possible questions thus are: Do we find similar patterns with regard to ordering in such languages? Is complexity a driving factor in these languages as well?

Another clear direction for future research only alluded to here is the precise structure of the left periphery. I suggested, in line with authors such as \citet{grohmann:2003:prolific}, \citet{laenzlinger:2005:french, laenzlinger:2011:elements} and \citet{kim:2019:syntax}, that the DP-external domain is universal. I showed that we need to assume this area in order to account for rather high direct and indirect modifiers, but also non-restrictive relative clauses, which originate there in Japanese. As for focused modifiers, they can productively appear in the left periphery; but they do not need to, and sometimes are forbidden from doing so. Future research should reveal the precise nature of this area. Questions that come to mind are: What other projections do we find in the left periphery? To what degree and according to which factors are the categories in this domain language-specific? Are there ordering principles in the left periphery, both from a cross-linguistic and language-specific perspective? \is{left periphery}

\subsection{Internal structure}

The final research desideratum concerns the internal structure of modifiers. The functional head \is{MOD@\textsc{mod}}\textsc{mod} as well as the complex structure of indirect modifiers proposed here are not only compatible with the universal DP structure, they also expand it significantly. Assuming the existence of \textsc{mod} in Japanese must mean this head exists in every language and languages differ in if and how they realize it. There are many similar elements which arguably are instances of \textsc{mod}, Persian \textit{E{\.z}āfe} or Mandarin \textsc{de} -- or generally all elements called \textit{linkers} \citep{richards:2023:finding} -- but of course in the majority of languages we do not see such elements. Even in Japanese, \textsc{no} is never present when overt tense is realized. How do we account for these observations? And can we draw generalizations from different languages that exhibit such elements?

I argue that the \is{MOD@\textsc{mod}}internal structure of modifiers should be viewed as important for research on the DP as the everlasting discussion on ordering restrictions and the (potential new) categories of functional projections within the DP-spine. Progress in this area will lead to an advance in research on the DP on a comparable scale and therefore deserves equal attention. In a nutshell, the internal structure of modifiers is another direction for future research on the DP.

\section{A final suggestion}

I will close this book with a suggestion: I think it is desirable that syntactic research keeps informing the field of Japanese Studies, or in fact any field of Area Studies, but importantly this also holds the other way around. The point is that these two disciplines can learn so much from one another, and so much potential remains untapped if researchers do not make use of it. Many Japanese Studies scholars will dismiss cartographic assumptions, also the ones defended here, immediately, as they seem too specific and rigid. Nevertheless, as I hope to have demonstrated, they make descriptive generalizations graspable by offering explananda thereof. On the other hand, the cartographic research field is dependent on specialists in the relevant fields, and only experts can grant them the deep understanding of relevant phenomena that is needed to derive the aforementioned explananda. My deep hope is to have inspired more studies in the future making use of both these research fields, Japanese studies and theoretical syntax, to enhance our knowledge in these two disciplines and beyond.

